\documentclass[../../main.tex]{subfiles}% 注意这里的文件路径不能用 ./main.tex ,否则用latexmk编译子文件会报错
\graphicspath{{\subfix{./image/}}} % 指定图片目录,后续可以直接使用图片文件名
% 注意这里的文件路径不能用 ../../image/ ,否则用latexmk编译子文件会报错

% 例如:
% \begin{figure}[H]
% \centering
% \includegraphics[scale=0.3]{图.png}
% \caption{}
% \label{figure:图}
% \end{figure}
% 注意:上述\label{}一定要放在\caption{}之后,否则引用图片序号会只会显示??.

\begin{document}

\section{分部积分公式与积分中值公式}

\begin{theorem}[分部积分公式]\label{theorem:实变函数--分部积分公式}
设 $f(x),g(x)$ 皆为 $[a,b]$ 上的可积函数, $\alpha,\beta\in\mathbb{R}$, 令
\[
F(x) = \alpha + \int_a^x f(t)\mathrm{d}t,\quad G(x) = \beta + \int_a^x g(t)\mathrm{d}t,
\]
则
\[
\int_a^b G(x)f(x)\mathrm{d}x + \int_a^b g(x)F(x)\mathrm{d}x = F(b)G(b) - F(a)G(a). \tag{5.13}
\]
\end{theorem}
\begin{remark}
注意到绝对连续函数与不定积分的关系, 上述定理可改述如下:

设 $f(x),g(x)$ 是 $[a,b]$ 上的绝对连续函数, 则
\[
\int_a^b f(x)g'(x)\mathrm{d}x + \int_a^b f'(x)g(x)\mathrm{d}x = f(b)g(b) - f(a)g(a).
\]
\end{remark}
\begin{proof}
设 $A,B$ 各为 $|F(x)|,|G(x)|$ 在 $[a,b]$ 上的最大值, 由
\[
|F(y)G(y)-F(x)G(x)| \leqslant A|G(y)-G(x)| + B|F(y)-F(x)|
\]
故知 $F(x)G(x)$ 是 $[a,b]$ 上的绝对连续函数, 从而 $F(x)G(x)$ 在 $[a,b]$ 上几乎处处可微, 且有
\[
(F(x)G(x))' = F(x)G'(x) + F'(x)G(x),\quad \text{a.e. } x\in[a,b].
\]
而对几乎处处的 $x\in[a,b]$, 有 $G'(x)=g(x)$ 与 $F'(x)=f(x)$, 故
\begin{align*}
\int_a^b G(x)f(x)\mathrm{d}x + \int_a^b F(x)g(x)\mathrm{d}x
&= \int_a^b G(x)F'(x)\mathrm{d}x + \int_a^b F(x)G'(x)\mathrm{d}x \\
&= \int_a^b (F(x)G(x))'\mathrm{d}x = F(b)G(b) - F(a)G(a).
\end{align*}

\end{proof}

\begin{example}
设 $f\in L([a,b])$, 且有
\[
\int_a^b x^n f(x)\mathrm{d}x = 0,\quad n=0,1,2,\cdots,
\]
则
\[
f(x)=0,\quad \text{a.e. } x\in[a,b].
\]
\end{example}
\begin{proof}
令 $F(x) = \int_a^x f(t)\mathrm{d}t$. 因为 $F(b)=0$, 且有
\begin{align*}
\int_a^b x^n f(x)\mathrm{d}x = x^n F(x)\Big|_a^b - n\int_a^b x^{n-1}F(x)\mathrm{d}x 
= -n\int_a^b x^{n-1}F(x)\mathrm{d}x = 0,\quad n=1,2,\cdots,
\end{align*}
所以我们得到
\[
\int_a^b x^n F(x)\mathrm{d}x = 0,\quad n=0,1,2,\cdots.
\]
现在, 根据\hyperref[Basis of Analytics-theorem:多项式逼近有限阶光滑函数]{多项式一致逼近连续函数定理}, 可知对任意的 $\varepsilon>0$, 存在多项式 $P(x)$, 使得 $|F(x)-P(x)|<\varepsilon$ ($x\in[a,b]$). 注意到
\[
\int_a^b P(x)F(x)\mathrm{d}x = 0,
\]
我们有
\[
\int_a^b F^2(x)\mathrm{d}x = \int_a^b F(x)(F(x)-P(x))\mathrm{d}x,
\]
从而可知
\[
\int_a^b F^2(x)\mathrm{d}x \leqslant \varepsilon\int_a^b |F(x)|\mathrm{d}x.
\]
由 $\varepsilon$ 的任意性可得 $F(x)\equiv 0$, 随之又得 $f(x)=0$, a.e. $x\in[a,b]$.

\end{proof}

\begin{theorem}[积分第一中值公式]\label{theorem:实变函数--积分第一中值公式}
若 $f(x)$ 是 $[a,b]$ 上的连续函数, $g(x)$ 是 $[a,b]$ 上的非负可积函数, 则存在 $\xi\in[a,b]$, 使得
\begin{align}
\int_a^b f(x)g(x)\mathrm{d}x = f(\xi)\int_a^b g(x)\mathrm{d}x. \label{eq::--2398ry89hjfjs5.14}
\end{align}
\end{theorem}
\begin{proof}
记 $f([a,b])=[c,d]$, 易知
\[
cg(x) \leqslant f(x)g(x) \leqslant dg(x),\quad \text{a.e. } x\in[a,b].
\]
将上式积分, 我们有
\[
c\int_a^b g(x)\mathrm{d}x \leqslant \int_a^b f(x)g(x)\mathrm{d}x \leqslant d\int_a^b g(x)\mathrm{d}x.
\]
若 $I = \int_a^b g(x)\mathrm{d}x > 0$, 则用它来除上式两端可得
\[
c \leqslant \frac{\int_a^b f(x)g(x)\mathrm{d}x}{I} \leqslant d.
\]
由 $f(x)$ 的连续性可知, 存在 $\xi\in[a,b]$, 使得
\[
f(\xi) = \frac{\int_a^b f(x)g(x)\mathrm{d}x}{\int_a^b g(x)\mathrm{d}x}.
\]
此即\eqref{eq::--2398ry89hjfjs5.14}式.

若 $I=0$, 则 $g(x)=0$, a.e. $x\in[a,b]$, 从而\eqref{eq::--2398ry89hjfjs5.14}式两端皆为零. 此时 $\xi$ 可任意地选取.

\end{proof}

\begin{theorem}[积分第二中值公式]\label{theorem:实变函数--积分第二中值公式}
若 $f\in L([a,b])$, $g(x)$ 是 $[a,b]$ 上的单调函数, 则存在 $\xi\in[a,b]$, 使得
\begin{align}
\int_a^b f(x)g(x)\mathrm{d}x = g(a)\int_a^\xi f(x)\mathrm{d}x + g(b)\int_\xi^b f(x)\mathrm{d}x. \label{eq::--2398ry89hjfjs5.15}
\end{align}
\end{theorem}
\begin{proof}
不妨设 $g(x)$ 是递增函数 (否则考查 $-g(x)$). 证明分两步进行:

(i) 假定 $g(x)$ 是 $[a,b]$ 上递增的绝对连续函数, 且令
\[
F(x) = \int_a^x f(t)\mathrm{d}t,\quad x\in[a,b].
\]
由分部积分公式可知
\[
\int_a^b f(x)g(x)\mathrm{d}x = F(b)g(b) - F(a)g(a) - \int_a^b F(x)g'(x)\mathrm{d}x.
\]
因为 $F(x)$ 是连续函数, 且 $g'(x)\geqslant 0$, a.e. $x\in[a,b]$, 所以由\eqref{eq::--2398ry89hjfjs5.14}式可知, 存在 $\xi\in[a,b]$, 使得
\[
\int_a^b F(x)g'(x)\mathrm{d}x = F(\xi)\int_a^b g'(x)\mathrm{d}x = F(\xi)(g(b)-g(a)).
\]
将上式代入前一式, 得到
\begin{align*}
\int_a^b f(x)g(x)\mathrm{d}x &= g(a)(F(\xi)-F(a)) + g(b)(F(b)-F(\xi)) \\
&= g(a)\int_a^\xi f(x)\mathrm{d}x + g(b)\int_\xi^b f(x)\mathrm{d}x.
\end{align*}

(ii) 对于 $g(x)$ 是递增的情形, 我们可以作一列递增且绝对连续的函数 $\{g_n(x)\}$, 使得 $g_n(a)=g(a)$, $g_n(b)=g(b)$ ($n=1,2,\cdots$), 且有
\[
\lim_{n\to\infty} g_n(x) = g(x),\quad \text{a.e. } x\in[a,b].
\]
$\{g_n(x)\}$的具体构造如下:

把 $[a,b]$ $n$ 等分, 令 $h_n = (b-a)/n$ 以及 $x_{n,k} = a + kh_n$, $0\leqslant k\leqslant n$, 作函数列:
\[
g_n(x) =
\begin{cases}
g(x), & x = x_{n,k}, \quad k=0,1,\cdots,n \quad (n\in\mathbb{N}), \\
\text{线性联结}, & x\in[x_{n,k-1},x_{n,k}], \quad k=1,2,\cdots,n \quad (n\in\mathbb{N}).
\end{cases}
\]
当 $x$ 是 $g$ 的连续点时有
\[
|g(x)-g_n(x)| \leqslant |g(x_{n,k-1})-g(x_{n,k})|,\quad x_{n,k-1}\leqslant x\leqslant x_{n,k}.
\]
对于 $g_n(x)$, 由 (i) 的结论, 则存在 $\xi_n\in[a,b]$, 使得
\[
\int_a^b f(x)g_n(x)\mathrm{d}x = g_n(a)\int_a^{\xi_n} f(x)\mathrm{d}x + g_n(b)\int_{\xi_n}^b f(x)\mathrm{d}x.
\]
这里不妨假定数列 $\{\xi_n\}$ 以 $\xi\in[a,b]$ 为极限 (否则可取子列), 于是令 $n\to\infty$, 上式转化为\eqref{eq::--2398ry89hjfjs5.15}式.

\end{proof}









\end{document}